\section{Vision Transformers and Segment-Based Visual Grouping}

Vision Transformers have become an important backbone family for image recognition. Unlike convolutional neural networks, which process images through local convolutional filters and spatial downsampling, Vision Transformers represent an image as a sequence of tokens and use self-attention to model relationships among them \cite{dosovitskiy2020image}. This token-based formulation provides a flexible framework for visual representation learning, but standard ViT models usually rely on fixed square patch tokens. Such tokens are convenient for computation, but because they are defined by a regular grid, they may cut across object boundaries and are not explicitly designed to correspond to semantic parts \cite{ke2024learning,park2025visually}. 

For fine-grained and hierarchical recognition, this limitation is important. Distinguishing bird species often requires localized evidence from object parts, such as the bill, wing, head, breast, or tail. At the same time, higher-level categories such as family or order may depend on broader visual structures. Therefore, a model that can organize visual information from small parts to larger regions is naturally suitable for hierarchical bird classification. This motivates the use of segment-based visual grouping methods in CAST \cite{ke2024learning} and H-CAST \cite{park2025visually}.


\subsection{Vision Transformers}

The Vision Transformer (ViT) adapts the Transformer architecture \cite{vaswani2017attention} to image classification by dividing an image into fixed-size patches and treating each patch as a token \cite{dosovitskiy2020image}. A learnable class token is appended to the patch tokens, and the sequence is processed by multiple Transformer encoder blocks. Through self-attention, each token can interact with other tokens, allowing the model to capture long-range dependencies across the image.

ViT provides a simple and powerful image classification framework, but its patch-based representation has limitations for visual parsing. Since patches are defined by a regular grid, they may cut across object boundaries or merge foreground and background pixels. This can be problematic in fine-grained recognition, where the relevant visual evidence may correspond to irregularly shaped parts rather than square image regions.

Another limitation is that standard ViT maintains patch tokens as the main visual units throughout the network. Although self-attention can aggregate information globally, the model does not explicitly produce a part-to-whole segmentation hierarchy. For hierarchical classification, this means that the model may not directly align fine-grained visual cues with fine labels and broader visual regions with coarse labels. These issues motivate later segment-based extensions of ViT, including CAST and H-CAST.


\subsection{CAST}

The model is named CAST because it Concurrently learns segmentation and recognition using Adaptive Segment Tokens, modifying the standard ViT formulation by replacing fixed square patches with adaptive segment tokens based on superpixels \cite{ke2024learning}. CAST enhances the standard ViT formulation with adaptive segment tokens and graph pooling, allowing segmentation and recognition to be learned concurrently within a single model \cite{ke2024learning}. Its central idea is that segmentation should not be treated only as an independent end task, but as an internal visual parsing process that evolves with and supports image recognition. CAST learns hierarchical segmentation together with recognition using image-level recognition objectives, without requiring pixel-level segmentation labels.

CAST starts by extracting superpixels from the input image using SEEDS \cite{van2015seeds}. These superpixels serve as the finest visual units and are used to construct adaptive segment tokens. Compared with fixed patch tokens, superpixels can better follow low-level color and boundary structure than fixed square patches, providing adaptive visual units before higher-level graph pooling. Each segment token is then processed by Transformer blocks, and graph pooling is used to progressively merge finer segments into coarser segments.

The graph pooling module constructs a fine-to-coarse hierarchy of segment tokens. At each level, finer segment tokens are softly assigned to coarser segment tokens based on feature similarity. This produces a hierarchy of segmentations and features, allowing the model to gradually move from local parts to larger object regions. The CAST paper emphasizes that this hierarchical segmentation is not trained with external segmentation labels; instead, it emerges from the recognition objective itself \cite{ke2024learning}.

This design distinguishes CAST from standard ViT in two major ways. First, the visual tokens are adaptive segments rather than fixed patches. Second, the number and granularity of tokens change through graph pooling, producing a part-to-whole visual hierarchy. As a result, CAST can provide both image-level recognition and internal hierarchical segmentation. The CAST paper reports that this segment-based representation improves unsupervised part-whole discovery, semantic segmentation, image classification, and efficiency compared with ViT baselines \cite{ke2024learning}.

CAST is relevant to this thesis because it provides the visual backbone on which H-CAST is built. Its fine-to-coarse segment hierarchy makes it possible to connect visual granularity with semantic granularity. In bird classification, fine segments can preserve local details such as the head, bill, or wing, while coarser segments can represent larger bird regions. This property is useful for hierarchical classification, where different taxonomy levels may require different levels of visual detail.


\subsection{H-CAST}

H-CAST extends CAST to hierarchical image classification \cite{park2025visually}. While CAST learns hierarchical segmentation for a flat recognition objective, H-CAST introduces hierarchical supervision so that different stages of the visual hierarchy are associated with different taxonomy levels. The model is motivated by the observation that fine-grained and coarse-grained classification may rely on different but related visual cues. For example, distinguishing between visually similar bird species may require local details such as the bill, wings, or head coloration, whereas recognizing a higher-level category may require broader object-level evidence.

The key idea of H-CAST is internal visual consistency. H-CAST argues that predictions at different hierarchy levels should be grounded in coherent and related visual regions. H-CAST addresses this issue by using the fine-to-coarse visual grouping structure inherited from CAST. Fine-level visual details are progressively grouped into coarser regions, so predictions at different hierarchy levels are encouraged to rely on coherent and related image regions.

Formally, H-CAST considers an image associated with labels at multiple hierarchy levels. It uses the segment and feature representations at different levels of the CAST hierarchy, adds a classification head at each level, and applies cross-entropy supervision across all hierarchy levels. In contrast to CAST, which uses a flat supervision signal on the final class token, H-CAST applies hierarchical supervision at multiple stages of the visual parsing process \cite{park2025visually}. This allows fine and coarse labels to guide different levels of visual representation.

In addition to visual consistency, H-CAST introduces Tree-Path KL Divergence loss to improve semantic consistency. This loss concatenates the labels from all hierarchy levels into a single distribution and aligns the concatenated outputs of the classification heads with this distribution using KL divergence. The purpose is to reduce prediction paths that violate the taxonomy, such as predicting a species whose parent family or order is inconsistent with the predicted higher-level labels \cite{park2025visually}.

The H-CAST paper evaluates hierarchical classification using both accuracy and consistency metrics. It defines Full-Path Accuracy (FPA) as the proportion of samples for which all hierarchy levels are correctly predicted, and uses Tree-based InConsistency Error rate (TICE) to measure whether predicted paths are valid in the taxonomy \cite{park2025visually}. These metrics are also relevant to this thesis because the proposed model is evaluated not only by species accuracy, but also by family accuracy, order accuracy, FPA, and TICE.

H-CAST is the main baseline and architectural foundation of this thesis. The proposed method keeps the H-CAST visual parsing pipeline and extends it with attribute-guided learning for CUB-200-2011. Specifically, instead of changing the hierarchical segment grouping mechanism, this thesis adds an attribute-guided module at the species-level branch. The motivation is that H-CAST provides visually consistent hierarchical representations, while CUB attributes provide additional semantic supervision about fine-grained bird traits. Combining these two sources of information may help the model better distinguish visually similar sibling species.